\chapter{容器、迭代器和标准库思维}

\section{问题入口：去重并保持顺序}

给定一组用户 ID，要求去掉重复项，但保留第一次出现的顺序。例如：

\begin{lstlisting}[language=C++,caption={输入输出示例}]
input  = {3, 5, 3, 7, 5, 9};
output = {3, 5, 7, 9};
\end{lstlisting}

朴素写法是每次检查结果数组里是否已有该元素：

\begin{lstlisting}[language=C++,caption={朴素去重}]
std::vector<int> unique_ids;
for (int id : ids) {
    bool seen = false;
    for (int existing : unique_ids) {
        if (existing == id) {
            seen = true;
            break;
        }
    }
    if (!seen) {
        unique_ids.push_back(id);
    }
}
\end{lstlisting}

这段代码正确，但复杂度是平方级。输入很大时，会反复扫描历史结果。

\section{组合容器表达需求}

需求有两个部分：判断是否出现过，保留顺序。哈希集合适合判断出现，向量适合保留顺序。

\begin{lstlisting}[language=C++,caption={哈希集合加顺序数组}]
#include <unordered_set>
#include <vector>

std::vector<int> deduplicate_keep_order(std::span<const int> ids)
{
    std::unordered_set<int> seen;
    std::vector<int> result;
    result.reserve(ids.size());

    for (int id : ids) {
        if (seen.insert(id).second) {
            result.push_back(id);
        }
    }
    return result;
}
\end{lstlisting}

\code{seen.insert(id).second} 表示这次是否真的插入了新元素。这个写法比先 \code{contains} 再 \code{insert} 少一次查找。

\section{vector 不是数组替身}

\code{std::vector} 是连续存储的动态数组。连续存储带来两个结果：遍历快，尾部追加通常快；但中间插入和删除需要移动后面的元素。

\begin{lstlisting}[language=C++,caption={删除元素时使用 erase-remove}]
std::vector<int> values{1, 2, 3, 2, 4};
const auto new_end = std::remove(values.begin(), values.end(), 2);
values.erase(new_end, values.end());
\end{lstlisting}

\code{std::remove} 不会真的缩短容器，它只是把保留元素移动到前面并返回新的逻辑结尾。真正删除尾部残余的是 \code{erase}。

\section{迭代器失效}

容器操作可能让迭代器、引用和指针失效。比如 \code{vector} 扩容后，原来指向元素的指针可能全都失效。

\begin{lstlisting}[language=C++,caption={尾部追加可能让指针失效}]
std::vector<int> values;
values.push_back(1);
int* first = &values[0];
values.push_back(2); // 可能扩容
// first 可能已经悬空
\end{lstlisting}

如果提前知道元素数量，用 \code{reserve} 可以减少扩容次数，但仍然不能把长期稳定地址建立在 \code{vector} 元素指针上。需要稳定节点地址时，考虑 \code{std::list}、\code{std::deque} 或自己管理对象存储，但要先确认真实需求。

\section{选择容器的经验表}

\begin{longtable}{p{0.24\linewidth}p{0.31\linewidth}p{0.35\linewidth}}
\toprule
需求 & 常用容器 & 注意点 \\
\midrule
顺序存储、频繁遍历 & \code{std::vector} & 扩容会移动元素 \\
两端插入删除 & \code{std::deque} & 不保证整体连续 \\
按 key 查找 & \code{std::unordered_map} & 平均快，依赖哈希质量 \\
有序遍历 & \code{std::map} & 节点结构，常数较大 \\
唯一集合 & \code{std::unordered_set} & 自定义类型需要 hash \\
\bottomrule
\end{longtable}

\begin{keyidea}
标准库容器不是语法点，而是复杂度合同。选择容器前先说清楚：是否需要顺序、是否需要唯一、是否需要稳定地址、主要操作是什么。
\end{keyidea}

\section{从输入规模推导复杂度}

朴素去重为什么慢？用一个具体输入推导。假设所有 ID 都不重复，且一共有 \code{n} 个。第一个元素比较 \code{0} 次，第二个比较 \code{1} 次，第三个比较 \code{2} 次，最后一个比较 \code{n - 1} 次。总比较次数是：

\[
0 + 1 + 2 + \cdots + (n - 1) = \frac{n(n - 1)}{2}
\]

这就是平方级。输入 \code{100} 个时还不明显，输入 \code{100000} 个时就会非常慢。哈希集合版本把“是否见过”变成平均常数时间查询，于是整体接近线性。算法优化不是凭感觉说“哈希快”，而是看历史信息是否被反复扫描。

\section{reserve 的作用和边界}

在去重函数里写 \code{result.reserve(ids.size())}，不是为了改变复杂度，而是减少扩容。因为最多保留 \code{ids.size()} 个元素，提前预留这个容量合理。注意 \code{reserve} 只改变容量，不改变元素个数：

\begin{lstlisting}[language=C++,caption={reserve 不会创建元素}]
std::vector<int> values;
values.reserve(10);
assert(values.size() == 0);
values.push_back(42);
assert(values.size() == 1);
\end{lstlisting}

如果写 \code{resize(10)}，就真的创建了十个元素。两者的区别必须分清：容量是能放多少，大小是已经有多少。

\section{迭代器失效要看容器合同}

不同容器的失效规则不同。入门阶段可以先抓住几个高频点：

\begin{itemize}
  \item \code{vector} 扩容会让所有元素引用、指针、迭代器失效。
  \item \code{vector} 中间删除会让删除点之后的位置失效。
  \item \code{list} 移动节点通常不让其他节点迭代器失效。
  \item \code{unordered_map} rehash 可能让迭代器失效，但元素引用通常仍指向元素对象。
\end{itemize}

不要凭容器名字猜。每次把迭代器保存到后面用，都要问：中间有没有可能插入、删除、扩容或 rehash？LRU 缓存选择 \code{std::list}，正是因为它需要在移动节点时保持迭代器稳定。

\section{容器选择案例：排行榜}

假设要维护游戏排行榜，需求是按玩家 ID 更新分数，并输出前十名。一个容器很难同时满足所有操作。可以先用 \code{unordered_map<int, int>} 保存玩家到分数的映射；输出前十时再把条目放进 \code{vector} 排序：

\begin{lstlisting}[language=C++,caption={排行榜的两阶段容器设计}]
using Scores = std::unordered_map<int, int>;

std::vector<std::pair<int, int>> top_scores(const Scores& scores)
{
    std::vector<std::pair<int, int>> items(scores.begin(), scores.end());
    std::ranges::sort(items, [](const auto& lhs, const auto& rhs) {
        if (lhs.second != rhs.second) {
            return lhs.second > rhs.second;
        }
        return lhs.first < rhs.first;
    });
    if (items.size() > 10) {
        items.resize(10);
    }
    return items;
}
\end{lstlisting}

如果更新非常频繁、查询前十也非常频繁，就可能需要额外维护有序结构。但第一版不要过早复杂化。先让主要操作和数据规模决定容器组合。

\begin{exercisebox}
给 \code{deduplicate_keep_order} 写三组测试：全重复、全不重复、混合重复。再把 \code{unordered_set} 去掉，手写朴素版本，对比两者在 \code{1000}、\code{10000}、\code{100000} 个随机 ID 下的运行时间。
\end{exercisebox}

\section{从需求字段推导容器组合}

容器选择不能只背“vector 快、map 有序”。先把需求字段列出来。去重并保持顺序这个问题有两条互相独立的需求：第一，能快速判断一个 ID 是否已经出现；第二，输出按第一次出现的顺序排列。单个容器很难同时把这两条都表达得最好。

\begin{longtable}{p{0.24\linewidth}p{0.28\linewidth}p{0.38\linewidth}}
\toprule
需求字段 & 合适结构 & 原因 \\
\midrule
快速判断出现过 & \code{std::unordered_set} & 平均常数时间查询和插入 \\
保留第一次出现顺序 & \code{std::vector} & 追加顺序就是输出顺序 \\
输出结果连续存放 & \code{std::vector} & 方便遍历、返回和测试 \\
\bottomrule
\end{longtable}

这样得到的解法不是“用了高级容器”，而是两个容器各自负责一个需求字段。以后做 LRU 缓存，会看到类似组合：哈希表负责按 key 找节点，链表负责维护最近使用顺序。

\section{地址稳定不是小细节}

如果代码保存了容器元素的地址，容器选择就进入另一个层次。看一个任务队列索引：

\begin{lstlisting}[language=C++,caption={保存 vector 元素指针的风险}]
std::vector<Task> tasks;
tasks.push_back(Task{.id = 1});
Task* first = &tasks[0];
tasks.push_back(Task{.id = 2});
first->done = true; // 如果扩容发生，first 可能已经失效
\end{lstlisting}

这不是 \code{vector} 的缺点，而是它的合同：为了保持连续存储，扩容时可以搬走所有元素。你要的是遍历快和紧凑内存，就接受地址可能变化；你要的是节点地址稳定，就选择别的结构或保存索引而不是指针。

一个实用替代是保存下标：

\begin{lstlisting}[language=C++,caption={保存下标而不是元素指针}]
std::size_t first_index = 0;
tasks.push_back(Task{.id = 2});
tasks[first_index].done = true;
\end{lstlisting}

下标也不是永远安全：如果删除或重排元素，它指向的逻辑任务可能变化。但它不会因为扩容变成悬空地址。容器合同要和修改操作一起看。

\section{rehash 也会改变迭代环境}

\code{unordered_map} 平均查找快，但插入很多元素时可能 rehash。rehash 会重新组织桶，迭代器会失效。假设你一边遍历一边插入：

\begin{lstlisting}[language=C++,caption={遍历哈希表时插入要小心}]
for (auto it = counts.begin(); it != counts.end(); ++it) {
    counts[new_key()] += 1; // 可能 rehash，使 it 失效
}
\end{lstlisting}

这类代码危险，因为循环变量本身可能被插入操作破坏。常见修复是先收集要插入的内容，遍历结束后统一插入；或者在已知规模时提前 \code{reserve}，并确认仍不违反容器规则。不要凭“我这次测试没崩”判断安全。

\section{erase 循环要让容器返回下一个位置}

删除元素时，迭代器推进规则也要严格。错误版本：

\begin{lstlisting}[language=C++,caption={删除后继续使用失效迭代器}]
for (auto it = values.begin(); it != values.end(); ++it) {
    if (*it < 0) {
        values.erase(it);
    }
}
\end{lstlisting}

\code{erase} 后 \code{it} 已经失效，循环末尾的 \code{++it} 不可靠。正确版本使用返回值：

\begin{lstlisting}[language=C++,caption={erase 返回下一个有效位置}]
for (auto it = values.begin(); it != values.end();) {
    if (*it < 0) {
        it = values.erase(it);
    } else {
        ++it;
    }
}
\end{lstlisting}

如果删除条件简单，优先使用 \code{std::erase_if}：

\begin{lstlisting}[language=C++,caption={用标准库表达删除条件}]
std::erase_if(values, [](int value) {
    return value < 0;
});
\end{lstlisting}

标准库算法不是为了炫技，而是把容易写错的迭代器协议封装成一个清晰意图。

\section{本章验收：写出容器账本}

选择容器前写出四列：主要操作、顺序要求、地址稳定要求、规模。比如“十万个用户，按 ID 查找，偶尔输出按注册顺序排列”，很可能需要 \code{unordered_map} 加 \code{vector} 或在对象里保存注册序号。比如“几百个规则，按优先级顺序扫描”，简单 \code{vector} 可能比树和哈希都好。容器选择的高级感不来自名字，而来自你能说清楚每个合同。
